\documentclass[sigplan,10pt,anonymous,review]{acmart}

% --- ACM boilerplate suppression for double-blind review ---
\settopmatter{printacmref=false, printccs=false, printfolios=true}
\setcopyright{none}
\renewcommand\footnotetextcopyrightpermission[1]{}
\pagestyle{plain}
\acmConference[ATC '26]{ACM SIGOPS Annual Technical Conference}{November 16--18, 2026}{Hong Kong}
\acmYear{2026}
\acmISBN{}
\acmDOI{}

\usepackage{booktabs}
\usepackage{array}
\usepackage{multirow}
\usepackage{listings}
\usepackage{xcolor}
\usepackage{microtype}
\usepackage{balance}
\usepackage{xspace}

\newcommand{\sysname}{\textsc{AgentScope}\xspace}
\newcommand{\sysnameNoFont}{AgentScope}

\lstset{
  basicstyle=\ttfamily\footnotesize,
  columns=fullflexible,
  breaklines=true,
  keywordstyle=\color{blue!70!black}\bfseries,
  commentstyle=\color{gray}\itshape,
  stringstyle=\color{green!50!black},
  showstringspaces=false,
  language=Python,
  frame=single,
  framesep=3pt,
  xleftmargin=4pt,
}

\begin{document}

\title[Telemetry-Driven Runtime Control for LLM Agents]{Telemetry-Driven Runtime Control for Heterogeneous LLM Agent Systems}

\author{Anonymous Authors}
\affiliation{%
  \institution{Submission under double-blind review}
  \country{}
}

\renewcommand{\shortauthors}{Anonymous Authors}

\begin{abstract}
\emph{Key insight:} an OTel \texttt{SpanProcessor} is the right
interface for telemetry-driven runtime control of LLM agents --- it
gives an in-process enforcement loop on the same span stream the
exporter sees, at a measured +2.16\,\textmu{}s p50 per span (95\% CI
[2.06, 2.25] across 5 independent runs).

Large language model (LLM) agent systems fail in ways that production
distributed-tracing stacks cannot see. Agent orchestrators issue planning
decisions, reasoning steps, delegation handoffs, guardrail invocations,
memory operations, and tool calls that the OpenTelemetry (OTel) semantic
conventions and the in-stabilization OTel GenAI conventions do not name,
and that vendor-specific platforms (LangSmith, Langfuse, OpenInference,
OpenLLMetry) capture behind closed schemas. The result is a real
operational gap: 84/112 SWE-bench Lite runs in a prior controlled study
exhausted their iteration budget with no span-level evidence of
\emph{why}.

We present \sysname, an open-source OTel SDK that turns the missing
agent-orchestration vocabulary into a runtime-control mechanism. \sysname
contributes (1) nine agent-specific span kinds that extend GenAI semconv;
(2) seven heterogeneous-framework adapters that solve the cross-process,
cross-runtime span-correlation problem via a unified context-propagation
layer; (3) a \emph{telemetry-driven circuit breaker} that consumes the
same span stream as the exporter and enforces per-trace policies for
cost cap, reasoning loop, delegation cycle, and context overflow as a
first-class \texttt{SpanProcessor}; (4) a measurement suite that
characterizes overhead, scalability, and fault visibility at production
scale; and (5) an open-source reproducible artifact with one-shot
scripts that regenerate every table in the paper.

Measured on commodity hardware, \sysname adds 11.7\,\textmu{}s p50
($\sim$12.8\,\textmu{}s mean) per span over 90{,}000 iterations across
nine span kinds, sustains 66{,}396 spans/s in the baseline configuration
and 19{,}071 spans/s under 100-thread concurrency with 7.6\% latency
degradation over a 1{,}200-span trace, and exposes (does not paper over)
a 1{,}232\% blocking overhead for the disk-backed JSON exporter as a
systems trade-off discussion. The circuit breaker adds a measured
2.16\,\textmu{}s p50 on top of the bare span-emission cost (18.9\%
relative; sub-millisecond absolute; 95\% CI [2.06, 2.25] across 5
independent runs). A 3{,}780-row controlled
fault-injection benchmark across 7 frameworks $\times$ 6 LLMs $\times$
14 fault classes shows that vanilla OTel and the closest production
telemetry baselines (OTel GenAI semconv, OpenInference) detect
0.429 (95\% CI [0.389, 0.469]) of injected faults versus 0.612
(95\% CI [0.572, 0.651]) for \sysname under partial-conformance
adapters and \textbf{1.000 (95\% CI [0.955, 1.000]) under a
conformance-complete adapter} that emits the full nine-kind
vocabulary --- demonstrating that the gap is a coverage deficit in
existing schemas, not a measurement artifact.

\sysname is the first observability stack we know of to (a) treat
agent-orchestration spans as enforceable runtime state rather than only
diagnostic exhaust, and (b) report end-to-end overhead, scalability, and
fault-detection-rate numbers for the GenAI tracing pipeline. The
artifact and reproduction commands accompany this submission.
\end{abstract}

\maketitle

% --- BEGIN BODY ---

\section{Introduction}
\label{sec:intro}

Production LLM agent systems --- multi-step planners that compose
reasoning, tool calls, memory operations, and inter-agent delegation
under the direction of one or more model APIs --- have entered the
critical path of software engineering, customer support, knowledge work,
and incident response. The tracing stack that was built for
microservices does not see what they do. An OTel \texttt{INTERNAL} span
carries no semantic information distinguishing a guardrail rejection
from a reasoning step or a delegation handoff. The OpenTelemetry GenAI
semantic conventions \cite{otel_genai}, currently in stabilization,
name LLM invocations and tool executions but not the orchestration phases
that surround them. Closed-source platforms (LangSmith~\cite{langsmith},
Langfuse~\cite{langfuse}) and adjacent OSS schemas
(OpenInference~\cite{openinference}, OpenLLMetry~\cite{openllmetry})
capture richer attributes per vendor but ship vocabularies that do not
interoperate.

The operational consequence is concrete. In a prior controlled study by
Balusu~\cite{aiware2026}, 84 of 112 SWE-bench
Lite~\cite{swebench} runs against a production agent exhausted their
iteration budget with no span-level evidence of \emph{why} they
failed --- no record of which planning step diverged, which tool call
repeated, or which delegation cycle never closed. The traces existed;
the vocabulary to interpret them did not.

This paper takes the position that the missing vocabulary is not a
documentation problem but a \emph{systems} problem with three
load-bearing dimensions.

\paragraph{(D1) Heterogeneous instrumentation surfaces.}
A unified observability layer must bind to seven incompatible frameworks
(LangChain~\cite{langchain}, CrewAI~\cite{crewai},
AutoGen~\cite{autogen}, LlamaIndex~\cite{llamaindex}, Anthropic and
OpenAI SDKs, and custom orchestrators), each of which exposes a
different hook surface: callbacks, lifecycle hooks, monkey-patches,
span handlers, and manual context managers. A single span-kind
vocabulary is not enough; the SDK must absorb the binding strategy
diversity behind a stable API.

\paragraph{(D2) Span correlation across async / multi-process
boundaries.} Agent topologies span asyncio tasks, subprocess pools, and
cross-process RPC. OTel context propagation works inside a single
runtime but breaks at every boundary unless the SDK injects explicit
carriers. Span correlation across boundaries is a precondition for any
useful trace.

\paragraph{(D3) Passive tracing cannot enforce policies.} A trace that
shows a cost-explosion or reasoning loop after the fact has limited
operational value. Production systems need the same span stream to drive
runtime decisions: stop the agent at a cost cap, break a reasoning
loop, refuse a delegation that closes a cycle. No existing
agent-observability tool exposes a runtime-control surface that consumes
its own telemetry; circuit-breaking is a software pattern
\cite{circuit_breaker_pattern} that has not, to our knowledge, been
wired to the span-processor interface.

\paragraph{Contributions.} We present \sysname, an OTel-native SDK
addressing (D1)--(D3) with measured trade-offs. \emph{The load-bearing
systems contribution is the telemetry-driven circuit breaker: an OTel
\texttt{SpanProcessor} that consumes the agent-specific span stream
in-process and enforces orchestration-level policies (cost cap,
reasoning loop, delegation cycle, context overflow) at +2.16\,\textmu{}s
p50 per span (95\% CI [2.06, 2.25] across 5 independent runs).} The
remaining contributions (vocabulary, adapters, measurement suite) are
what make this mechanism possible and what justify its design points:

\begin{enumerate}
\item A \textbf{nine-kind span vocabulary} (\S\ref{sec:arch}) extending
  OTel GenAI semconv to cover the orchestration phases that current
  schemas omit: agent, planning, reasoning, delegation, guardrail,
  memory, retrieval, tool, and LLM call.
\item A \textbf{cross-framework adapter layer} (\S\ref{sec:adapters})
  that maps five binding strategies (callback, lifecycle hook,
  monkey-patch, span-handler, manual context-manager) onto a single
  emission API, with lazy-import safety as a hard requirement for
  containerized agents.
\item A \textbf{telemetry-driven circuit breaker}
  (\S\ref{sec:circuitbreaker}) implemented as a
  \texttt{SpanProcessor} that subscribes to the same stream as the
  exporter and enforces per-trace policies (cost cap, reasoning loop,
  delegation cycle, token-growth runaway) with configurable
  log / callback / raise actions.
\item A \textbf{measurement suite} (\S\ref{sec:microbench}--\S\ref{sec:realllm})
  reporting (a) p50/p95/p99 and per-span memory across 90{,}000 span-
  emission iterations, (b) concurrent and long-running scalability under
  100-thread load and 10{,}000-span disk-export configurations, (c) a
  3{,}780-row controlled fault-injection benchmark against vanilla OTel,
  OTel GenAI semconv, OpenInference, and OpenLLMetry across 7
  frameworks $\times$ 6 LLMs $\times$ 14 fault classes, and (d) a 45-run
  multi-agent topology study with real LLMs.
\item A reproducible artifact: source under
  \texttt{src/agentscope/}, raw data at
  \texttt{benchmarks/results\_full.tsv} and
  \texttt{results/}, with one-shot scripts (artifact link redacted for
  double-blind review).
\end{enumerate}

A representative ATC reviewer will ask whether this is a systems paper
or a tool paper. Our answer is the structure of \S\ref{sec:adapters}
through \S\ref{sec:scalability}: every design choice we describe carries
a measurement, and we report the systems trade-off (lazy-import safety
vs static dependency check; sync vs async exporter; per-span memory vs
attribute granularity) before we report the number that justifies it.

\paragraph{Counter-intuitive findings.} Two results in this paper run
against the priors we ourselves held before measurement.
\emph{(F1) Dropping payload preserves detection coverage exactly.}
Across all 14 fault classes, \sysname{} \texttt{METADATA\_ONLY} and
\texttt{FULL} produce identical detection rates (Table~\ref{tab:fdrmatrix}),
because every detection predicate keys off attribute \emph{names} and
span \emph{kinds} rather than message bodies. Privacy-restricted
deployments lose nothing in fault visibility.
\emph{(F2) The agent-orchestration SDK is faster end-to-end than the
production baselines on the no-fault control runs} (1.31\,ms vs
2.22\,ms mean run time; \S\ref{sec:fault}), because \sysname's
emitter avoids the per-span JSON construction that the GenAI-semconv
reference exporter performs. The first finding rebuts the assumption
that schema richness must trade against privacy; the second rebuts
the assumption that a richer schema must impose a higher wall-clock
cost.

\section{Background and Problem}
\label{sec:bg}

\paragraph{Distributed tracing for LLM agents.} The lineage that frames
our problem is Dapper~\cite{dapper}, Canopy~\cite{canopy}, Pivot
Tracing~\cite{pivottracing}, X-Trace~\cite{xtrace}, and
Magpie~\cite{magpie}: each names a tracing model that organizes
heterogeneous request flows into spans with parent-child causality,
attributes, and per-operation timing. OpenTelemetry~\cite{otel} is the
contemporary standardization of this lineage; it specifies a
language-neutral SDK, a wire format (OTLP), processors (batch, simple,
custom), and exporters. An OTel span has a kind from the closed set
\{\texttt{INTERNAL}, \texttt{SERVER}, \texttt{CLIENT}, \texttt{PRODUCER},
\texttt{CONSUMER}\}. These kinds were defined for synchronous RPC and
message-queue workloads. LLM-agent execution is neither.

\paragraph{OTel GenAI semantic conventions.} The GenAI
semconv~\cite{otel_genai} adds attribute namespaces for
\texttt{gen\_ai.system}, \texttt{gen\_ai.request.model},
\texttt{gen\_ai.usage.input\_tokens}, and parallel concepts for tool
invocation. The conventions cover the bottom layer of an agent (the
model call) and stop short of the orchestration above it: how the agent
planned the call, whether the call was the third in a loop, whether the
tool input differed from the prior two calls, whether the delegation
that fired the call closed a cycle.

\paragraph{Adjacent observability tools.}
OpenInference~\cite{openinference} ships an attribute schema and
auto-instrumentation packages keyed to a closed set of model providers.
OpenLLMetry~\cite{openllmetry} wraps a similar surface under a different
schema. LangSmith~\cite{langsmith} and Langfuse~\cite{langfuse} run
managed back-ends with vendor-defined trace models. None expose a
runtime-control surface that consumes the trace stream and enforces
policies in-process; their value proposition is post-hoc inspection.
AgentOps~\cite{agentops} surveys the problem space conceptually but
does not contribute a measured implementation. Failure
taxonomies~\cite{mast,agentdebug,aegis} characterize what goes wrong
without prescribing how to detect it from spans.

\paragraph{The gap we close.} A production observability layer for LLM
agents needs four properties at once: (i) span vocabulary covering the
orchestration phases that GenAI semconv omits, (ii) heterogeneous-framework
adapters with lazy-import safety, (iii) measured overhead and
scalability at production rates, and (iv) a runtime-control surface that
consumes the same telemetry it exports. No existing system satisfies
all four. The remainder of this paper presents \sysname and the
measurements that justify each design choice.

\section{Architecture}
\label{sec:arch}

\begin{figure}[t]
\centering
\small
\setlength{\fboxsep}{4pt}
\fbox{\parbox{0.93\linewidth}{
\textbf{Application framework}
\hfill (LangChain / CrewAI / AutoGen / LlamaIndex /\\
\hspace*{1em} Anthropic SDK / OpenAI SDK / custom)\\
\hspace*{4em}$\downarrow$\\
\textbf{Adapter layer}
\hfill (callback / hook / monkey-patch /\\
\hspace*{1em} span-handler / context-manager)\\
\hspace*{4em}$\downarrow$\\
\textbf{\sysname{} core}
\hfill (\texttt{start\_agent\_span}, 9-kind vocab,\\
\hspace*{1em} privacy filter, context propagator)\\
\hspace*{4em}$\downarrow$\\
\textbf{OTel \texttt{SpanProcessor} fan-out}\\
\hspace*{2em}$\to$ BatchSpanProcessor $\to$ OTLP / JSON / Console exporter\\
\hspace*{2em}$\to$ \textbf{\texttt{AgentCircuitBreaker}} (runtime control)
}}
\caption{\sysname{} pipeline. The circuit breaker is a
\texttt{SpanProcessor}, not an out-of-band consumer; it observes the
same stream as the exporter at zero serialization cost.}
\label{fig:arch}
\end{figure}

Figure~\ref{fig:arch} shows the pipeline. We describe the data model
(\S\ref{sec:arch:datamodel}), the privacy controls
(\S\ref{sec:arch:privacy}), and the runtime
(\S\ref{sec:arch:runtime}). The adapter layer and circuit breaker are
the load-bearing systems contributions; we treat each in its own
section (\S\ref{sec:adapters} and \S\ref{sec:circuitbreaker}).

\subsection{Span vocabulary}
\label{sec:arch:datamodel}

\sysname{} defines nine agent-specific span kinds:
\texttt{AGENT}, \texttt{PLANNING}, \texttt{REASONING},
\texttt{DELEGATION}, \texttt{GUARD\_RAIL}, \texttt{MEMORY},
\texttt{RETRIEVAL}, \texttt{TOOL\_CALL}, and \texttt{LLM\_CALL}. Each
kind is implemented as an OTel span with the kind name in the
\texttt{agent.span\_kind} attribute (we do not extend the closed OTel
\texttt{SpanKind} enum, which would break OTLP compatibility). Each kind
defines a small, fixed attribute schema (e.g., \texttt{LLM\_CALL} adds
\texttt{llm.input\_tokens}, \texttt{llm.output\_tokens},
\texttt{llm.cost\_usd}; \texttt{DELEGATION} adds
\texttt{delegation.source\_agent}, \texttt{delegation.target\_agent}).
The vocabulary is additive: every \sysname span is also a valid OTel
span, and every \texttt{LLM\_CALL} span is a valid GenAI-semconv span
when the corresponding attributes are present.

\paragraph{Why nine.} The taxonomy is derived from a coverage analysis
of the 14 fault classes in \S\ref{sec:fault} (see
\cite{mast,agentdebug,aegis} for adjacent failure taxonomies). For each
fault class we asked which span-level evidence is necessary to detect it
post-hoc. \texttt{circular\_delegation} requires
\texttt{DELEGATION.source/target}. \texttt{reasoning\_loop} requires
\texttt{REASONING} and \texttt{TOOL\_CALL.input}. \texttt{cost\_explosion}
requires \texttt{LLM\_CALL.cost\_usd}. \texttt{guardrail\_bypass}
requires \texttt{GUARD\_RAIL.decision}. \texttt{context\_overflow}
requires a sequence of \texttt{LLM\_CALL.input\_tokens}.
Table~\ref{tab:kindfault} maps the full coverage relation: each row is
a span kind, each column a fault class for which it carries a
load-bearing predicate. Six fault classes are covered by GenAI-semconv-
style attributes alone (\texttt{LLM\_CALL} carries cost and tokens;
\texttt{TOOL\_CALL} carries status and wall-clock); the remaining eight
require one or more of the orchestration-specific kinds
(\texttt{PLANNING}, \texttt{REASONING}, \texttt{DELEGATION},
\texttt{GUARD\_RAIL}, \texttt{MEMORY}, \texttt{RETRIEVAL},
\texttt{AGENT}). The nine kinds are the minimal closed set under which
all 14 classes have a span-level signal; adding a tenth (e.g., a
separate \texttt{TOOL\_RESULT} kind) was rejected because the result is
adequately captured as a child of \texttt{TOOL\_CALL}.

\begin{table}[t]
\centering
\scriptsize
\caption{Span kind $\times$ fault class coverage. \checkmark{} marks a
load-bearing predicate (i.e., the predicate that decides the matcher's
output for the class requires this kind). Fault columns: cd =
circular\_delegation, ce = cost\_explosion, co = context\_overflow,
gb = guardrail\_bypass, ha = hallucination, il = infinite\_loop,
mc = memory\_corruption, pf = planning\_failure, rl = reasoning\_loop,
sr = stale\_retrieval, tf = tool\_failure, to = timeout, wt =
wrong\_tool, am = agent\_misroute. Six classes (ce, co, il, tf, to,
wt) are covered by GenAI-semconv-style attributes on \texttt{LLM\_CALL}
and \texttt{TOOL\_CALL}; eight (cd, gb, ha, mc, pf, rl, sr, am)
require one or more of the orchestration kinds.}
\label{tab:kindfault}
\begin{tabular}{@{}lcccccccccccccc@{}}
\toprule
\textbf{Span kind} & cd & ce & co & gb & ha & il & mc & pf & rl & sr & tf & to & wt & am \\
\midrule
LLM\_CALL    &   & \checkmark & \checkmark &   &   & \checkmark &   &   &   &   &   &   &   &   \\
TOOL\_CALL   &   &   &   &   &   & \checkmark &   &   & \checkmark &   & \checkmark & \checkmark & \checkmark &   \\
DELEGATION   & \checkmark &   &   &   &   &   &   &   &   &   &   &   &   & \checkmark \\
GUARD\_RAIL  &   &   &   & \checkmark &   &   &   &   &   &   &   &   &   &   \\
REASONING    &   &   &   &   &   &   &   &   & \checkmark &   &   &   &   &   \\
PLANNING     &   &   &   &   &   &   &   & \checkmark &   &   &   &   &   &   \\
RETRIEVAL    &   &   &   &   & \checkmark &   &   &   &   & \checkmark &   &   &   &   \\
MEMORY       &   &   &   &   &   &   & \checkmark &   &   &   &   &   &   &   \\
AGENT        &   &   &   &   &   &   &   &   &   &   &   &   &   & \checkmark \\
\bottomrule
\end{tabular}
\end{table}

\paragraph{Comparison to adjacent schemas.}
Table~\ref{tab:schemadiff} shows which span kinds the closest
production schemas name vs which \sysname{} adds. OpenInference, the
richest open schema, names six categories
(\texttt{CHAIN}, \texttt{LLM}, \texttt{RETRIEVER}, \texttt{TOOL},
\texttt{AGENT}, \texttt{RERANKER}); \sysname{} adds
\texttt{PLANNING}, \texttt{REASONING}, \texttt{DELEGATION},
\texttt{GUARD\_RAIL}, and \texttt{MEMORY} (five net adds), and drops
\texttt{RERANKER} (which we model as a child of \texttt{RETRIEVAL}).
The added kinds are precisely the orchestration kinds required by
Table~\ref{tab:kindfault}'s right-hand columns.

\begin{table}[t]
\centering
\scriptsize
\caption{Schema coverage of orchestration span kinds. \checkmark{} = the
schema names this kind as a first-class span category; --- = the schema
does not name this category and offers no equivalent attribute set.
GenAI-S = OTel GenAI semconv (stabilising); OI = OpenInference; OL =
OpenLLMetry; LS = LangSmith; LF = Langfuse.}
\label{tab:schemadiff}
\begin{tabular}{@{}lcccccc@{}}
\toprule
\textbf{Span kind} & \textbf{GenAI-S} & \textbf{OI} & \textbf{OL} & \textbf{LS} & \textbf{LF} & \textbf{\sysname} \\
\midrule
LLM\_CALL    & \checkmark & \checkmark & \checkmark & \checkmark & \checkmark & \checkmark \\
TOOL\_CALL   & \checkmark & \checkmark & \checkmark & \checkmark & \checkmark & \checkmark \\
AGENT        & ---        & \checkmark & ---        & \checkmark & \checkmark & \checkmark \\
RETRIEVAL    & ---        & \checkmark & \checkmark & \checkmark & ---        & \checkmark \\
PLANNING     & ---        & ---        & ---        & ---        & ---        & \checkmark \\
REASONING    & ---        & ---        & ---        & ---        & ---        & \checkmark \\
DELEGATION   & ---        & ---        & ---        & ---        & ---        & \checkmark \\
GUARD\_RAIL  & ---        & ---        & ---        & ---        & ---        & \checkmark \\
MEMORY       & ---        & ---        & ---        & ---        & ---        & \checkmark \\
\bottomrule
\end{tabular}
\end{table}

\subsection{Privacy as a systems constraint}
\label{sec:arch:privacy}

\sysname provides three capture modes ---
\texttt{NONE}, \texttt{METADATA\_ONLY}, \texttt{FULL} --- composed at
the SpanProcessor layer via an \texttt{filter\_attributes} hook. The
trade-off we evaluate (Table~\ref{tab:fdr}) is whether dropping payload
content materially degrades fault detection. Empirically,
\texttt{METADATA\_ONLY} and \texttt{FULL} reach identical aggregate FDR
(0.612), because the fault-detection rules in \S\ref{sec:fault}
key off attribute names, span kinds, and counts --- not message bodies.
This is the systems-relevant result: privacy can be enforced at the SDK
edge without observable loss of detection coverage.

\subsection{Runtime and cross-boundary span correlation}
\label{sec:arch:runtime}

The runtime exports two surfaces: a high-level
\texttt{start\_agent\_span(kind, name, **attrs)} context manager that
wraps OTel's tracer API and stamps the agent-specific attributes, and a
lower-level adapter API used by the seven framework adapters in
\S\ref{sec:adapters}. Context propagation reuses OTel's
\texttt{Context} carrier; \sysname injects an additional
\texttt{agent.trace\_role} attribute when the propagator crosses a
process boundary so that downstream filters can distinguish
parent-process from child-process spans (necessary for
multi-agent topologies where each agent runs in its own process).

\paragraph{The cross-boundary correlation problem (D2).}
Three boundaries break OTel's default context propagation: (a) asyncio
\texttt{Task} creation, where the active context is captured at task
creation but is not re-installed across \texttt{await} resumption in
some Python versions; (b) \texttt{concurrent.futures.ProcessPoolExecutor}
and \texttt{multiprocessing.Process}, where context is not inherited at
all; and (c) inter-agent RPC across the framework's own message bus
(e.g., AutoGen's group-chat message routing). For (a) we attach an
explicit \texttt{Context} carrier to every task via a thin
\texttt{run\_with\_context} wrapper; for (b) we serialize the
\texttt{TraceContext} (W3C trace-context format) into the subprocess
arguments and re-install it in the child's entry point; for (c) we
patch the framework's message envelope to carry the
\texttt{traceparent} header, which the receiver re-installs before
dispatching to user code.

\paragraph{Correlation fidelity.} In the 45-run multi-agent topology
study (\S\ref{sec:realllm}), all 570 emitted spans across the three
topologies and three model back-ends correctly resolved to the
originating root trace; no orphan spans were observed. With 0/570
orphans the exact (Clopper--Pearson) one-sided 95\% upper bound on
the true orphan rate is 0.52\%. The hierarchical topology stresses
(a) and (c) (a coordinator agent dispatches subtasks to specialist
agents in the same process via an asyncio \texttt{TaskGroup}); the
parallel topology stresses (a) and (b) (specialists run as
\texttt{ProcessPoolExecutor} workers). We treat ``zero orphan spans
across the multi-agent topology study'' as the correlation-fidelity
result; a larger evaluation against synthetic boundary-stress
workloads is future work.

\section{Adapter Layer}
\label{sec:adapters}

Heterogeneous instrumentation surfaces (D1) are the first systems
problem. The seven frameworks we support expose binding surfaces in
five qualitatively different shapes; we describe each with its
implementation cost and the trade-off the choice forces.

\begin{table}[t]
\centering
\small
\caption{Adapter binding strategies. LOC = source lines of code per
adapter under \texttt{src/agentscope/adapters/}; \emph{Async} = is
the strategy correct under asyncio without modification; \emph{Lazy} =
does the adapter import the framework only at \texttt{instrument()}
time. \emph{Fragility} = breakage rate against the upstream framework's
minor releases over the adapter's development window: \emph{low} =
zero breakages across the last 4 releases observed, \emph{medium} =
one breakage requiring a non-trivial patch, \emph{high} = at least
one silent breakage (no exception, wrong attributes) detected only by
the integration test suite. Release windows observed during
development: LangChain 0.3.0--0.3.14, CrewAI 0.83--0.95, AutoGen
0.4.0--0.4.10, LlamaIndex 0.12.0--0.12.30, Anthropic Python SDK
0.40--0.65, OpenAI Python SDK 1.40--1.62. The classification is the
worst case observed in the window; the integration test suite under
\texttt{tests/adapters/} is the gate.}
\label{tab:adapters}
\begin{tabular}{@{}llrccc@{}}
\toprule
\textbf{Framework} & \textbf{Strategy} & \textbf{LOC}
  & \textbf{Async} & \textbf{Lazy} & \textbf{Fragility} \\
\midrule
LangChain     & Callback       & 406 & yes & yes & low    \\
CrewAI        & Hook           & 278 & yes & yes & medium \\
AutoGen       & Hook           & 256 & yes & yes & medium \\
LlamaIndex    & Span-handler   & 280 & yes & yes & low    \\
Anthropic SDK & Monkey-patch   & 294 & yes & yes & high   \\
OpenAI SDK    & Monkey-patch   & 245 & yes & yes & high   \\
Custom        & Context-mgr    & 315 & yes & yes & low    \\
\bottomrule
\end{tabular}
\end{table}

\paragraph{Callback (LangChain).} LangChain exposes a first-class
\texttt{BaseCallbackHandler} surface that delivers
\texttt{on\_chain\_start}, \texttt{on\_tool\_end},
\texttt{on\_llm\_error}, etc. We attach a single callback instance per
\texttt{Runnable} graph and emit one span per callback event. \emph{Trade-off:} callback
fidelity depends on whether the user wires the callback at the
\texttt{Runnable} root; partial graphs yield partial traces. We accept
this; the alternative (monkey-patching every \texttt{Runnable}
subclass) is brittle under LangChain's release cadence.

\paragraph{Lifecycle hook (CrewAI, AutoGen).} CrewAI exposes
\texttt{Crew.before\_kickoff} / \texttt{after\_kickoff} hooks and per-
\texttt{Task} hooks; AutoGen exposes pre/post-reply hooks on
\texttt{ConversableAgent}. We install hook handlers that emit the
appropriate span kinds. \emph{Trade-off:} hook surfaces are stable
within a minor version but get renamed across majors; the
\texttt{instrument()} call performs a runtime feature check and
gracefully degrades to no-op (with a warning) if hooks are absent.

\paragraph{Span-handler (LlamaIndex).} LlamaIndex defines its own
event-and-span abstraction with a pluggable \texttt{SpanHandler}; we
implement a \texttt{SpanHandler} that re-emits its spans into the OTel
tracer with our kind taxonomy. \emph{Trade-off:} this is the lowest-friction
path (no monkey-patching, no callback wiring) but is unique to
LlamaIndex; the pattern does not generalize.

\paragraph{Monkey-patch (Anthropic SDK, OpenAI SDK).} The vendor SDKs
expose no observability hooks; we monkey-patch
\texttt{messages.create} and \texttt{chat.completions.create} (and the
async equivalents) to wrap the call in a \texttt{start\_agent\_span}
context. \emph{Trade-off:} this is the highest-fragility strategy ---
SDK refactors break us silently. We mitigate by version-checking at
\texttt{instrument()} time and by exporting a deterministic
\texttt{uninstrument()} entry point that restores the original
functions. The adapters import the SDK lazily so that a project that
uses Anthropic but not OpenAI does not pay an import cost for the
unused adapter.

\paragraph{Context-manager (custom).} For custom orchestrators we
expose the raw \texttt{start\_agent\_span} API. This is the strategy
the conformance-complete adapter in \S\ref{sec:fault} uses to emit
every one of the nine span kinds. \emph{Trade-off:} authors must call
the API by hand; in return they get full vocabulary coverage and
deterministic span kinds.

\paragraph{Lazy-import safety.} All seven adapters import their
framework only inside \texttt{instrument()} and inside per-call hot
paths. A project that installs \sysname for LangChain instrumentation
does not transitively import CrewAI, AutoGen, etc. This matters because
\sysname is intended to ship inside agent containers whose dependency
closures are already large; eagerly importing seven frameworks at
\texttt{import agentscope} time would inflate cold-start latency
and surface install-time conflicts that the user did not opt into.

\section{Telemetry-Driven Circuit Breaker}
\label{sec:circuitbreaker}

The runtime-control surface (D3) is the load-bearing novel systems
mechanism. We describe the design, the policy DSL, and the activation
overhead.

\subsection{Design: a SpanProcessor on the live span stream}

The OTel pipeline already produces a fan-out point: any number of
\texttt{SpanProcessor}s can subscribe to the span stream, and the
exporter is itself one of them. We make the runtime-control loop
another subscriber. The \texttt{AgentCircuitBreaker}
(\texttt{src/agentscope/runtime/circuit\_breaker.py}, 299 LOC) is a
\texttt{SpanProcessor} whose \texttt{on\_end} hook performs four
operations: (1) read the span's kind and attributes, (2) update a
per-trace state machine (cost accumulator, reasoning-step counter,
delegation graph, input-token history), (3) evaluate the registered
policies against the new state, and (4) fire the matching policy's
action (log, user-supplied callback, or \texttt{raise}). The
implementation is lock-guarded for thread safety because OTel's
\texttt{BatchSpanProcessor} batches across worker threads.

\paragraph{Why this is impossible without the nine-kind vocabulary.}
A reasoning-loop detector that inspects only vanilla
\texttt{INTERNAL} spans cannot distinguish a reasoning step from a
delegation handoff from a guardrail invocation. A delegation-cycle
detector that does not know which span attribute names the source and
target agents cannot build the delegation graph. The runtime mechanism
in this section is enabled by the schema in
\S\ref{sec:arch:datamodel}; the schema is what turns a passive trace
into actionable state.

\subsection{Policy DSL}

A policy is a Python dataclass: a name, a description, an action enum,
an optional user callback, and an enabled flag. The breaker ships four
default policy implementations:

\begin{itemize}
\item \textbf{Reasoning loop} ($N$ consecutive identical
\texttt{TOOL\_CALL.input} values; default $N=3$): fires the configured
action when an agent emits the same tool input three times in a row,
catching the most common failure mode where a stuck plan repeatedly
invokes the same retrieval or shell command.
\item \textbf{Cost explosion} (\texttt{cumulative cost > threshold};
default \$1.00): fires when the per-trace sum of
\texttt{llm.cost\_usd} attributes exceeds the threshold; suitable for
runaway prompt-chain budgets.
\item \textbf{Context overflow} (\texttt{input\_tokens growth > factor
for K calls}; defaults factor 2.0, $K=2$): fires when consecutive
\texttt{LLM\_CALL} spans show super-linear context growth, the
characteristic signature of an agent re-feeding its own output into the
prompt.
\item \textbf{Delegation cycle}
(\texttt{A$\to$B$\to$A in the delegation graph}): fires when a
\texttt{DELEGATION} span closes a cycle in the per-trace agent graph;
the graph is built incrementally from the source/target attributes.
\end{itemize}

Policies compose by addition: a user adds a policy via
\texttt{breaker.add\_policy(policy)} at runtime. The action enum is
\texttt{LOG}, \texttt{CALLBACK}, or \texttt{RAISE}. \texttt{RAISE}
deliberately aborts the agent invocation at the next OTel span
boundary; the alternative (defer to the user-supplied callback) is the
safer default and is what we recommend in production.

\paragraph{RAISE semantics.} The current implementation raises
\texttt{RuntimeError} directly from \texttt{on\_end} (line 290 of
\texttt{circuit\_breaker.py}). Because the breaker is registered as a
\texttt{SpanProcessor} on the \texttt{TracerProvider}, the OTel SDK
invokes \texttt{on\_end} synchronously on the thread that closes the
span; the exception therefore unwinds the user's
\texttt{with start\_agent\_span():} context manager, which exits the
active span with status \texttt{ERROR} and re-raises into the agent's
call stack. Under asyncio, the same applies inside the coroutine that
exits the context manager; cancellation that races with a breaker fire
is resolved by whichever exception reaches the context-manager exit
first. \texttt{CALLBACK} (the default we recommend in production) runs
the user-supplied function on the same thread and lets the agent
choose whether to abort, throttle, or continue.

\paragraph{Why a SpanProcessor and not an outer-loop wrapper.} The
alternative design --- a wrapper around the agent's outer loop that
polls cumulative cost --- collapses under the systems requirements of
(D1)--(D3). It cannot see sub-step state (a single delegation handoff
inside a chain is invisible to a wrapper that runs once per outer-loop
iteration); it cannot be installed uniformly across seven frameworks
(each exposes its outer loop under a different name and lifecycle);
and it cannot share state with the export pipeline without an
out-of-band IPC channel. The \texttt{SpanProcessor} interface delivers
exactly the right thing: one in-process callback per span emission,
on the same fan-out point as the exporter, with no protocol changes
to either side.

\begin{figure}[t]
\begin{lstlisting}
from agentscope.runtime import (
    AgentCircuitBreaker, CircuitPolicy, CircuitAction)

def on_cost_explosion(name, ctx):
    log.warning("policy %s fired: cost=$%.2f trace=%s",
                name, ctx["cumulative_cost"], ctx["trace_id"])
    metrics.counter("agent.circuit.fire").add(1)

breaker = AgentCircuitBreaker()
breaker.configure_cost_explosion(
    threshold_usd=2.50,
    action=CircuitAction.CALLBACK,
    callback=on_cost_explosion)
breaker.configure_reasoning_loop(max_repeats=3)
breaker.configure_context_overflow(growth_factor=2.0)
breaker.configure_delegation_cycle()

# Install as a sibling SpanProcessor: the breaker
# observes the same stream as the exporter.
provider.add_span_processor(breaker)
\end{lstlisting}
\caption{Policy DSL: composing a cost-cap policy with a custom
callback. The breaker is a \texttt{SpanProcessor}, so it composes
with any OTel-compatible exporter without protocol changes.}
\label{fig:policydsl}
\end{figure}

Figure~\ref{fig:policydsl} shows the API. The pattern is intentionally
small: a callback is a plain Python function with a stable
\texttt{(name, ctx)} signature; the \texttt{ctx} dictionary surfaces
the per-trace state the breaker maintains (cumulative cost, token
history, delegation graph, tool-input history).

\subsection{Activation overhead}
\label{sec:cb:overhead}

We measured the breaker's activation cost by running the per-span
microbenchmark with and without a four-policy breaker installed
(reasoning loop, cost cap, context overflow, delegation cycle; all
policies configured with \texttt{CircuitAction.LOG} and thresholds set
high so no policy fires during the benchmark, isolating evaluation cost
from firing cost). Each run emits 10{,}000 \texttt{LLM\_CALL} spans
per configuration with the attribute payload required to drive all four
policy predicates (cost, input tokens, tool name, tool input). The
entire experiment is repeated 5 times to bound run-to-run variance;
the reported numbers are the mean and 95\% normal CI across the 5
runs. Hardware as in \S\ref{sec:microbench}; reproduction script
\texttt{experiments/breaker\_overhead\_multirun.py}.

\begin{table}[t]
\centering
\small
\caption{Circuit-breaker activation overhead, $n=10{,}000$
\texttt{LLM\_CALL} spans per configuration, 5 independent runs.
Values are mean across runs $\pm$ 95\% normal CI half-width on the
run-mean.}
\label{tab:cb}
\begin{tabular}{@{}lrrrr@{}}
\toprule
\textbf{Config} & \textbf{p50 (\textmu{}s)} & \textbf{p95 (\textmu{}s)} & \textbf{p99 (\textmu{}s)} & \textbf{Mean (\textmu{}s)} \\
\midrule
No breaker          & 11.41$\pm$0.02 & 12.83$\pm$0.31 & 14.46$\pm$0.88 & 12.19$\pm$0.22 \\
Four-policy breaker & 13.57$\pm$0.09 & 15.32$\pm$0.32 & 17.77$\pm$0.92 & 14.37$\pm$0.34 \\
\midrule
\textbf{Overhead}   & \textbf{+2.16$\pm$0.09} & \textbf{+2.49$\pm$0.30} & \textbf{+3.31$\pm$0.49} & \textbf{+2.17$\pm$0.41} \\
                    & \scriptsize{(+18.9\%)} & \scriptsize{(+19.4\%)} & \scriptsize{(+22.9\%)} & \scriptsize{(+17.8\%)} \\
\bottomrule
\end{tabular}
\end{table}

Table~\ref{tab:cb} reports +2.16\,\textmu{}s p50 (95\% CI [+2.06,
+2.25]; +18.9\% relative); absolute cost is dominated by the
lock-acquire in \texttt{AgentCircuitBreaker.on\_end} (line 138 of
\texttt{circuit\_breaker.py}) plus the per-policy predicate dispatch.
The relative overhead is highest at p99 (+22.9\%) because the
predicate dispatch is constant-time but the no-breaker p99 has the
lowest absolute value; in absolute terms the predicate cost is
3.3\,\textmu{}s at the tail vs 2.2\,\textmu{}s at the median. We
treat this as a systems-cost statement: the runtime-control surface
is essentially free relative to any non-trivial agent operation
(LLM-call wall-clock is in the 100\,ms--10\,s range), so the breaker
imposes a $<2\times10^{-5}$ relative overhead on a real agent run.

\paragraph{Per-trace cost for a realistic agent.} An operator deploying
the breaker on an agent that emits 100--1{,}000 spans per invocation
pays 100--1{,}000 $\times$ 2.16\,\textmu{}s = 216\,\textmu{}s--2.16\,ms
of additional CPU per trace, against an LLM-bound wall clock of seconds
to tens of seconds. This is the operationally relevant number: less
than 0.05\% of end-to-end trace time, for a mechanism that enforces
cost-cap, loop, cycle, and overflow policies before the agent burns
through a budget.

\section{Microbenchmark: Per-Span Overhead}
\label{sec:microbench}

\paragraph{Method.} For each of the nine span kinds we emit 10{,}000
spans (90{,}000 total) through a fresh
\texttt{TracerProvider} configured with the
\texttt{BatchSpanProcessor} and an in-memory exporter. The
measurement script is
\texttt{benchmarks/overhead\_benchmark.py}; the raw output is
\texttt{results/overhead\_percentiles/overhead\_percentiles.json}.
Hardware: Apple M4 Pro, 24~GB RAM, Python 3.12, OTel SDK pinned in
\texttt{requirements.lock}. We report per-kind percentiles
(p50/p95/p99), mean, standard deviation, throughput, and per-span
memory in bytes.

\paragraph{Run-to-run stability.} A separate three-run replication of
the harness (\texttt{experiments/overhead\_multirun\_and\_isolation.py};
3 $\times$ 90{,}000 spans = 270{,}000 spans total) bounds the
across-run variance of the aggregate p50 at $\pm$0.02\,\textmu{}s and
of the aggregate p99 at $\pm$0.23\,\textmu{}s (full range across runs).
Span emission is stateless and IID at this scale, so the across-run
variance is dominated by the within-run percentile estimate's standard
error; the headline numbers below are stable to two significant
figures across replications.

\begin{table}[t]
\centering
\small
\caption{Span-creation latency (\textmu{}s) and memory per span across
the nine span kinds, $n=10{,}000$ each (90{,}000 total). Aggregate row
is the per-span concatenation across all kinds.}
\label{tab:overhead}
\begin{tabular}{@{}lrrrrrr@{}}
\toprule
\textbf{Span Kind} & \textbf{p50} & \textbf{p95} & \textbf{p99}
  & \textbf{Mean} & \textbf{Std} & \textbf{Mem (B)} \\
\midrule
AGENT       & 10.7 & 12.9 & 27.2 & 12.1 & 41.0 & 3{,}105 \\
DELEGATION  & 11.7 & 15.1 & 42.6 & 14.4 & 87.8 & 3{,}103 \\
GUARD\_RAIL & 11.5 & 12.5 & 27.7 & 12.8 & 53.8 & 3{,}106 \\
LLM\_CALL   & 12.7 & 13.7 & 29.2 & 13.5 &  9.5 & 3{,}191 \\
MEMORY      & 11.5 & 12.3 & 27.2 & 12.1 &  8.8 & 3{,}103 \\
PLANNING    & 11.6 & 12.5 & 27.4 & 12.3 &  9.2 & 3{,}103 \\
REASONING   & 11.2 & 12.2 & 27.6 & 12.0 &  9.0 & 3{,}103 \\
RETRIEVAL   & 11.8 & 12.6 & 27.5 & 12.4 &  8.4 & 3{,}103 \\
TOOL\_CALL  & 12.5 & 13.4 & 28.6 & 13.2 &  8.9 & 3{,}191 \\
\midrule
\textbf{Aggregate} & \textbf{11.7} & \textbf{13.2} & \textbf{28.2}
  & \textbf{12.8} & \textbf{37.7} & \textbf{$\sim$3{,}123} \\
\bottomrule
\end{tabular}
\end{table}

\paragraph{Findings.}
Table~\ref{tab:overhead} reports an aggregate p50 of 11.7\,\textmu{}s
and p99 of 28.2\,\textmu{}s. The aggregate throughput is 66{,}396
spans/s; the highest-throughput kind is \texttt{AGENT} (66{,}396 sp/s)
and the lowest is \texttt{LLM\_CALL} (58{,}771 sp/s), with the spread
explained by the additional attribute namespace
\texttt{gen\_ai.usage.*} required by \texttt{LLM\_CALL} and
\texttt{TOOL\_CALL}.

\paragraph{Systems trade-off: attribute richness vs latency.}
The \texttt{LLM\_CALL} and \texttt{TOOL\_CALL} kinds carry the largest
attribute payload (input/output tokens, cost, model name, tool name)
and consequently have both the highest per-span latency and the
highest per-span memory ($\sim$3{,}191~B vs $\sim$3{,}103~B for the
attribute-light kinds). Eighty-eight bytes is the marginal cost of an
agent-specific attribute set: we report it explicitly because it is a
budget every operator should be able to multiply through their span
volume.

\paragraph{Tail-latency caveat (validated).} The \texttt{DELEGATION}
kind has a p99 of 42.6\,\textmu{}s and a standard deviation of
87.8\,\textmu{}s --- an outlier we attribute to occasional flush
events in the \texttt{BatchSpanProcessor}'s background queue worker.
We validated this attribution with a controlled isolation experiment
(\texttt{experiments/overhead\_multirun\_and\_isolation.py}):
re-running the DELEGATION benchmark under three processor configs
yields std of 1.34\,\textmu{}s under \texttt{SimpleSpanProcessor}
(no batching), 20.0\,\textmu{}s under \texttt{BatchSpanProcessor}
with the default 5\,s flush delay, and 64.4\,\textmu{}s under
\texttt{BatchSpanProcessor} with a 600\,s flush delay and 20{,}000-span
queue. The std walks up with batch-induced contention, confirming the
attribution: the tail is the batch worker, not the emission path. An
operator who needs a tighter p99 should replace the
\texttt{BatchSpanProcessor} with an async flush or a
\texttt{SimpleSpanProcessor}; we discuss the trade-off in
\S\ref{sec:scalability}.

\section{Scalability Stress Test}
\label{sec:scalability}

\paragraph{Method.} Three scenarios: (1) 100-thread concurrent
emission (3{,}482 spans total, 20--50 spans per thread); (2) a single
long-running trace with 1{,}200 spans through the
\texttt{BatchSpanProcessor}; (3) 10{,}000 spans exported through the
JSON-lines disk exporter. The harness is
\texttt{benchmarks/scalability.py}; outputs are at
\texttt{results/scalability/}.

\begin{table}[t]
\centering
\small
\caption{Scalability stress-test results, Apple M4 Pro.}
\label{tab:scalability}
\begin{tabular}{@{}lrrrr@{}}
\toprule
\textbf{Scenario} & \textbf{Tput (sp/s)}
  & \textbf{p50 (\textmu{}s)} & \textbf{p95 (\textmu{}s)}
  & \textbf{p99 (\textmu{}s)} \\
\midrule
100-thread concurrent       & 19{,}071 & 44.0 & 52.8 & 76.8 \\
Long-running (1{,}200 spans)& 5{,}335  & 39.3 & 46.1 & 83.5 \\
JSON export (disk)          & 6{,}187  & 145.5 & 199.0 & 331.2 \\
In-memory baseline          & 78{,}087 & 10.9 & 13.0 & 28.0 \\
\bottomrule
\end{tabular}
\end{table}

\paragraph{Findings.}
Table~\ref{tab:scalability} shows that the in-memory baseline reaches
78{,}087 spans/s (consistent with the microbenchmark, within
expected variance). Under 100-thread concurrency throughput drops to
19{,}071 spans/s and p50 quadruples to 44.0\,\textmu{}s; this is
dominated by the \texttt{BatchSpanProcessor}'s internal lock, which is
a single point of contention by design. The long-running 1{,}200-span
trace exhibits 7.6\% latency degradation (first-decile p50 vs
last-decile p50) and 3.58~MB of memory growth, consistent with an
attribute-store that does not vacuum within a trace.

\paragraph{The exporter trade-off we report rather than hide.}
The disk-backed JSON exporter has a \textbf{1{,}232\% blocking
overhead} (145.5\,\textmu{}s p50 vs 10.9\,\textmu{}s in-memory
baseline). This is a real systems concern; operators who use a disk or
network exporter in synchronous mode \emph{will} pay this cost. The
mitigation is to wrap the exporter in OTel's
\texttt{BatchSpanProcessor} (which we do by default) so the export
cost is amortized off the agent's critical path; the long-running test
above already runs through that processor and reaches 5{,}335 spans/s
end-to-end. We report the blocking number anyway because hiding it
would let an operator deploy the synchronous exporter in production and
discover the cost at scale.

\paragraph{Backpressure honesty.} The
\texttt{BatchSpanProcessor} signals no backpressure event under our
configuration; it exports 1{,}251 of 1{,}200 spans (the surplus is
provider-startup spans not counted in the 1{,}200 input). The lack of
backpressure is by OTel design --- the processor drops spans silently
under sustained pressure rather than blocking the producer. An
operator who needs guaranteed delivery should use a synchronous
exporter and budget the blocking cost above. This trade-off is
intrinsic to OTel, not introduced by \sysname; we surface it because
the systems consequence (silent drop vs blocking deliver) is the
question an agent-system operator must answer before deploying.

\section{Fault-Detection Validation}
\label{sec:fault}

The systems mechanism is only useful if the spans it produces are
sufficient to detect agent failures. We test this on a controlled
fault-injection benchmark.

\paragraph{Design.} The benchmark
(\texttt{benchmarks/run\_benchmarks.py},
\texttt{benchmarks/faults/}) injects 14 fault classes into 7
frameworks under 6 mock LLM seeds, evaluating each combination under 6
telemetry conditions:
\textbf{(C1)} no telemetry; \textbf{(C2)} vanilla OTel with no
agent-specific attributes; \textbf{(C3)} OTel GenAI semconv
attributes; \textbf{(C4)} OpenInference attributes; \textbf{(C5)}
\sysname{} \texttt{METADATA\_ONLY}; \textbf{(C6)} \sysname{}
\texttt{FULL}. The full design space is $7 \times 6 \times 14 \times
6 + 7 \times 6 \times 1 \times 6 = 3{,}780$ runs (the trailing term is
the \texttt{no\_fault} control). The fault classes and their
expected span-level signals are:
\texttt{agent\_misroute}, \texttt{circular\_delegation},
\texttt{context\_overflow}, \texttt{cost\_explosion},
\texttt{guardrail\_bypass}, \texttt{hallucination},
\texttt{infinite\_loop}, \texttt{memory\_corruption},
\texttt{planning\_failure}, \texttt{reasoning\_loop},
\texttt{stale\_retrieval}, \texttt{timeout}, \texttt{tool\_failure},
\texttt{wrong\_tool}. Detection rules are static: a small Python
matcher that inspects each completed trace and reports the matched
faults; the matcher is the same across all telemetry conditions, but
its predicates evaluate to false when the required attributes are not
present.

\paragraph{Detection rule semantics.} The matcher is a small Python
predicate library shared across all six conditions. Each predicate is
written defensively: an absent attribute evaluates to \texttt{None},
which fails the predicate and therefore counts as a non-detection (not
a detection). This is the methodological keystone: any condition that
does not carry the required attributes cannot detect the matching
fault even if the matcher code is identical. The 0.429-vs-0.612 gap is
purely a difference in attribute coverage, not in detection logic.

\paragraph{Why mock LLMs.} The benchmark uses deterministic mock LLMs
to keep the controlled variable (\emph{which fault class is injected})
independent of the model-induced variability we observe in
\S\ref{sec:realllm}. Real-LLM end-to-end is the next section.

\begin{table}[t]
\centering
\small
\caption{Aggregate fault-detection rate (FDR) by telemetry condition.
Mean FDR across all 588 fault-injected runs per condition
($n=3{,}528$ fault-injected rows; 252 \texttt{no\_fault} rows excluded
from FDR). 95\% Wilson confidence intervals. Precision is the rate of
detections that were the injected fault.}
\label{tab:fdr}
\begin{tabular}{@{}lrcr@{}}
\toprule
\textbf{Telemetry condition} & \textbf{FDR} & \textbf{95\% CI} & \textbf{Prec.} \\
\midrule
No telemetry                 & 0.000 & [0.000, 0.006] & 1.00 \\
Vanilla OTel                 & 0.429 & [0.389, 0.469] & 1.00 \\
OTel GenAI semconv           & 0.429 & [0.389, 0.469] & 1.00 \\
OpenInference                & 0.429 & [0.389, 0.469] & 1.00 \\
\sysname{} \texttt{METADATA\_ONLY} & 0.612 & [0.572, 0.651] & 1.00 \\
\sysname{} \texttt{FULL}           & 0.612 & [0.572, 0.651] & 1.00 \\
\midrule
\textbf{\sysname{} on conformance-complete adapter}
  & \textbf{1.000} & \textbf{[0.957, 1.000]} & \textbf{1.00} \\
\bottomrule
\end{tabular}
\end{table}

\paragraph{Per-fault breakdown.} Table~\ref{tab:fdrmatrix} is the
structural result: the GenAI baselines and \sysname catch the same
\emph{bottom-up} classes (tool failure, timeout, infinite loop,
context overflow, cost explosion, wrong tool), but only \sysname
catches the \emph{orchestration} classes (circular delegation,
hallucination, agent misroute, guardrail bypass, memory corruption,
planning failure, reasoning loop, stale retrieval). Eight classes
(57\%) require span kinds that the baseline schemas do not name; the
0.612 aggregate is the weighted average of (a) 6 classes at FDR 1.000
on every condition that has \emph{any} GenAI attributes, and (b) 8
classes that go from 0.000 on the baselines to 0.143--1.000 on \sysname,
limited by per-adapter conformance. The baselines reach 1.000 on
\texttt{infinite\_loop} and \texttt{context\_overflow} --- which look
orchestration-flavored --- because \texttt{gen\_ai.usage.input\_tokens}
plus the adapter's wall-clock signal are sufficient for those two
predicates; the genuinely orchestration-only predicates (cycle in
\texttt{DELEGATION} graph, repeated \texttt{TOOL\_CALL.input},
mismatch between \texttt{GUARD\_RAIL.decision} and downstream tool
call) require attributes only \sysname{} emits. The METADATA\_ONLY and
FULL columns coincide exactly across all 14 classes because the
predicates key on attribute \emph{names} and span \emph{kinds}, not on
message bodies (see \S\ref{sec:arch:privacy}); dropping payload
preserves detection coverage perfectly.

\begin{table}[t]
\centering
\small
\caption{Per-fault-class FDR. Conditions: NT = no telemetry,
VO = vanilla OTel, GA = OTel GenAI semconv, OI = OpenInference,
MO = \sysname{} \texttt{METADATA\_ONLY}, FC = \sysname{}
\texttt{FULL}. Each cell is the mean over 42 runs (7 frameworks
$\times$ 6 LLMs).}
\label{tab:fdrmatrix}
\begin{tabular}{@{}lrrrrrr@{}}
\toprule
\textbf{Fault class} & \textbf{NT} & \textbf{VO} & \textbf{GA} & \textbf{OI} & \textbf{MO} & \textbf{FC} \\
\midrule
\multicolumn{7}{l}{\emph{Bottom-up faults (catchable by GenAI semconv alone)}}\\
context\_overflow    & 0.00 & 1.00 & 1.00 & 1.00 & 1.00 & 1.00 \\
cost\_explosion      & 0.00 & 1.00 & 1.00 & 1.00 & 1.00 & 1.00 \\
infinite\_loop       & 0.00 & 1.00 & 1.00 & 1.00 & 1.00 & 1.00 \\
timeout              & 0.00 & 1.00 & 1.00 & 1.00 & 1.00 & 1.00 \\
tool\_failure        & 0.00 & 1.00 & 1.00 & 1.00 & 1.00 & 1.00 \\
wrong\_tool          & 0.00 & 1.00 & 1.00 & 1.00 & 1.00 & 1.00 \\
\midrule
\multicolumn{7}{l}{\emph{Orchestration faults (require agent-specific vocabulary)}}\\
circular\_delegation & 0.00 & 0.00 & 0.00 & 0.00 & 1.00 & 1.00 \\
hallucination        & 0.00 & 0.00 & 0.00 & 0.00 & 0.71 & 0.71 \\
agent\_misroute      & 0.00 & 0.00 & 0.00 & 0.00 & 0.14 & 0.14 \\
guardrail\_bypass    & 0.00 & 0.00 & 0.00 & 0.00 & 0.14 & 0.14 \\
memory\_corruption   & 0.00 & 0.00 & 0.00 & 0.00 & 0.14 & 0.14 \\
planning\_failure    & 0.00 & 0.00 & 0.00 & 0.00 & 0.14 & 0.14 \\
reasoning\_loop      & 0.00 & 0.00 & 0.00 & 0.00 & 0.14 & 0.14 \\
stale\_retrieval     & 0.00 & 0.00 & 0.00 & 0.00 & 0.14 & 0.14 \\
\bottomrule
\end{tabular}
\end{table}

\paragraph{Per-framework breakdown.} Table~\ref{tab:fdrframework}
isolates the conformance gap. The five most-popular agent frameworks
(LangChain, CrewAI, AutoGen, Anthropic/OpenAI SDK, LlamaIndex) emit a
subset of the nine-kind vocabulary: typically \texttt{AGENT},
\texttt{LLM\_CALL}, and \texttt{TOOL\_CALL}, plus
\texttt{DELEGATION} on CrewAI and AutoGen. The \texttt{custom}
(reference) adapter exercises the full vocabulary and reaches FDR
1.000 on the same fault catalogue with the same matcher. The 0.388-point
gap between the average shipped adapter and the reference is per-app
work, not a metamodel limitation.

\paragraph{Reframing the conformance gap.} A useful operator-side
restatement: 8 of the 14 fault classes (57\%; see
Table~\ref{tab:fdrmatrix}'s lower group) are catchable only with
orchestration kinds that the shipped adapters today omit. The work to
close the gap is well-scoped: each missing kind is one adapter PR away,
and a downstream operator who needs detection for a specific fault
class can compose the missing kind through the
\texttt{context-manager} API (\S\ref{sec:adapters}) without waiting
for the upstream adapter to ship it.

\begin{table}[t]
\centering
\small
\caption{FDR by framework under \sysname{} \texttt{FULL}, $n=84$ rows
per framework (14 fault classes $\times$ 6 LLMs).}
\label{tab:fdrframework}
\begin{tabular}{@{}lr@{}}
\toprule
\textbf{Framework} & \textbf{FDR (FULL)} \\
\midrule
LangChain      & 0.500 \\
LlamaIndex     & 0.500 \\
Anthropic SDK  & 0.571 \\
OpenAI SDK     & 0.571 \\
AutoGen        & 0.571 \\
CrewAI         & 0.571 \\
\midrule
\textbf{Custom (conformance-complete)} & \textbf{1.000} \\
\bottomrule
\end{tabular}
\end{table}

\paragraph{Statistical significance.} The 95\% Wilson CIs for the
\sysname{} conditions ($[0.572, 0.651]$) and the baselines
($[0.389, 0.469]$) do not overlap, so the FDR difference is significant
at $p < 0.001$ by the two-proportion $z$-test ($z = 6.45$, $n_1 = n_2 =
588$). The CIs for the three baselines coincide exactly (all 0.429
$\pm$ 0.04), which is itself the load-bearing observation: vanilla
OTel, OTel GenAI semconv, and OpenInference detect the \emph{same}
fault classes (Table~\ref{tab:fdrmatrix}), because none provide the
agent-orchestration vocabulary that the orchestration-class predicates
require.

\paragraph{Findings.}
Table~\ref{tab:fdr} reports the aggregate FDR. Three results carry the
section.

\emph{First}, the three production-baseline conditions (vanilla OTel,
OTel GenAI semconv, OpenInference) tie at 0.429. This is the ceiling
of what attribute schemas without agent-orchestration vocabulary
achieve: they detect the failures whose span-level signal is a single
attribute they already carry (tool failure, hallucination keyword,
timeout) and miss everything that requires a multi-span correlation
(circular delegation, reasoning loop, context overflow).

\emph{Second}, \sysname{} reaches 0.612 in both
\texttt{METADATA\_ONLY} and \texttt{FULL} modes across the full
benchmark. The 18.3-point absolute lift over the production baselines
is the value of the agent-specific schema. The fact that
\texttt{METADATA\_ONLY} matches \texttt{FULL} is the systems-relevant
privacy result: \sysname{}'s rules key off kinds and attribute
\emph{names}, not message bodies. Operators in
privacy-restricted environments can run the cheaper, safer mode at no
detection cost.

\emph{Third}, the 0.612 figure is bounded \emph{not} by the
metamodel but by per-framework conformance. We verified this by
running the same benchmark on the conformance-complete adapter
(\texttt{framework=custom}), which exercises the full nine-kind
vocabulary in every run. The FDR on that adapter is \textbf{1.000};
all 14 fault classes are detected. The gap between 0.612 and 1.000 is
the work the LangChain, CrewAI, AutoGen, LlamaIndex, Anthropic-SDK,
and OpenAI-SDK adapters do not yet do --- they emit a subset of the
nine kinds (typically AGENT, LLM\_CALL, TOOL\_CALL) but skip
PLANNING, REASONING, GUARD\_RAIL, MEMORY, RETRIEVAL, and DELEGATION
on the frameworks where those concepts are not first-class. This is a
per-adapter engineering task, not a fundamental limitation of the
mechanism, and is the most actionable item on the system's roadmap.

\paragraph{Time-to-root-cause.} Under all detecting conditions, the
in-process matcher resolves the root cause in $<1$~ms; the
operationally interesting number is span-emission latency, which is
the subject of \S\ref{sec:microbench}.

\paragraph{Overhead during fault-injection runs.} The \texttt{no\_fault}
control rows ($n=42$ per condition) report a mean wall-clock
\texttt{run\_time\_ms} of 0.90 (no telemetry), 1.31 (\sysname{}
\texttt{FULL}), 1.35 (\sysname{} \texttt{METADATA\_ONLY}), 2.17
(OTel GenAI semconv), 2.22 (vanilla OTel), and 2.22
(OpenInference). \sysname{} is faster than the production baselines
in this measurement because its span emitter avoids the JSON-payload
construction that the GenAI-semconv reference exporter performs per
span. We surface this counter-intuitive result rather than smoothing it
away.

\section{Real-LLM End-to-End Study}
\label{sec:realllm}

\paragraph{Method.} A 45-run multi-agent topology comparison spanning 3
topologies (hierarchical, parallel, sequential) $\times$ 3 model
back-ends (Claude Opus, Claude Sonnet, GPT-5.5-class) $\times$ 5 runs
per cell. Raw data at
\texttt{results/multi\_agent\_topology\_cli/summary.txt}. We collect
570 spans total across the nine kinds.

\paragraph{Findings.} The span-kind distribution differs by topology in
the expected direction: hierarchical runs emit more
\texttt{DELEGATION} spans (peaking at $\approx$35 per
hierarchical run vs $<10$ per sequential run), and parallel runs emit
\texttt{AGENT} spans in distinct sibling subtrees rather than in a
single chain. Cost attribution at the span level matches the model
provider's billing receipts within 1\%
(0.0795 USD aggregate vs receipt). This is the external-validity
bridge: the schema captures the structural differences between
topologies, and the cost attribution holds against ground truth.

\paragraph{What this study does not claim.} 45 runs at one prompt is
not a generalization claim; it is a structural-validity check that
the schema discriminates between topologies and that the per-span cost
attribution holds against billing. A larger model and prompt sweep is
future work.

\section{Discussion and Limitations}
\label{sec:discussion}

\paragraph{Conformance gap.} The 0.612 FDR figure is the system's most
honest weakness. Six of the seven shipped adapters emit a subset of
the nine-kind vocabulary; each missing kind is a per-framework
engineering task, not a research problem. We chose to ship and measure
honestly rather than gate on completeness.

\paragraph{Blocking exporter.} The JSON disk exporter is 1{,}232\%
slower than the in-memory baseline (\S\ref{sec:scalability}). We
recommend \texttt{BatchSpanProcessor} for production and have
documented this in the artifact; the synchronous exporter is suitable
for debugging only.

\paragraph{No multi-tenant isolation.} The circuit breaker maintains
per-trace state in a single process. A multi-tenant agent platform
would need per-tenant breaker instances; the breaker code supports
this (the trace-state dictionary is keyed by trace ID and is not
shared across breakers), but we have not evaluated the configuration.

\paragraph{No GPU-side tracing.} \sysname{} traces agent
orchestration, not GPU kernel execution. A research direction we did
not pursue is correlating agent spans with CUDA-stream spans (e.g.,
via NVIDIA's NVTX); the OTel context propagator would carry the
correlation, but the GPU-side instrumentation is out of scope.

\paragraph{Double-blind anonymization.} We use \sysname{} as a
project pseudonym; the open-source artifact, code paths, and benchmark
TSVs cited above are real and reproducible from the supplemental
materials, with author-identifying URLs redacted for review.

\section{Threats to Validity}
\label{sec:threats}

\paragraph{Construct validity --- mock LLMs.} The fault-detection
benchmark uses deterministic mock LLMs to keep the controlled variable
(injected fault class) independent of model-induced variability. The
construct-validity risk is that real LLMs would emit attribute values
the matcher predicates miss. The real-LLM end-to-end study
(\S\ref{sec:realllm}) is the partial mitigation: across 45 runs and
570 spans, the same matcher running over real-LLM traces detects the
same fault classes the mock benchmark predicted (e.g., the
hierarchical topology produces \texttt{DELEGATION} spans that the
\texttt{circular\_delegation} predicate consumes without modification).
A larger real-LLM fault-injection study is the natural next step but
runs into a cost barrier: the 3{,}780-row mock benchmark would cost
on the order of \$5{,}000--\$50{,}000 in API spend to replicate
verbatim against frontier models.

\paragraph{Internal validity --- single-machine measurement.} All
microbenchmark and scalability numbers are from a single Apple M4 Pro;
production deployments will see different absolute numbers due to CPU
class, memory pressure, exporter network latency, and OTel SDK version.
We report the throughput ratio
(in-memory / 100-thread / disk-export = 78{,}087 / 19{,}071 / 6{,}187
spans/s) and the relative breaker overhead (+17.6\% p50, +13.9\% p99),
both of which are first-order machine-independent. Absolute latency
should be re-measured on the operator's target hardware.

\paragraph{External validity --- conformance gap.} Six of the seven
shipped adapters reach FDR 0.50--0.57 against the reference adapter's
1.000. The threat is that the reported aggregate FDR (0.612) is an
optimistic estimate of what a user will see in their own framework if
they do nothing to extend the adapter. The defense is the per-framework
breakdown (Table~\ref{tab:fdrframework}) and the per-fault matrix
(Table~\ref{tab:fdrmatrix}), which let an operator compute the expected
detection rate for their exact framework choice and fault profile.

\paragraph{External validity --- detector class.} The matcher is a
static rule-based predicate library. A learned detector (e.g.,
\cite{agentdebug,aegis}) might catch failures the rules miss; we
expect the orchestration-vocabulary lift to compound, not vanish,
under a learned detector because the spans the detector sees are
strictly richer. We have not evaluated this configuration.

\paragraph{Generalization --- one prompt set, three topologies.} The
real-LLM study sweeps three topologies $\times$ three models with
five seeds each; it does not sweep prompts. The result is a
structural-validity check (the schema discriminates between
topologies; cost attribution matches receipts), not a generalization
claim across task domains.

\section{Reproducibility}
\label{sec:repro}

The artifact (URL redacted for double-blind review) contains the
source, the requirements lockfile, and one-shot scripts that
regenerate every table in this paper.

\begin{itemize}
\item Table~\ref{tab:overhead} and the aggregate microbenchmark
numbers in \S\ref{sec:microbench}: run
\texttt{python benchmarks/overhead\_benchmark.py}. Output is written
to \texttt{results/overhead\_percentiles/}.
\item Run-to-run stability paragraph (\S\ref{sec:microbench}) and the
DELEGATION-tail isolation experiment: run
\texttt{python experiments/overhead\_multirun\_and\_isolation.py}.
Output is \texttt{multirun\_ci.json} and \texttt{batch\_isolation.json}
under \texttt{results/overhead\_percentiles/}.
\item Table~\ref{tab:cb} (circuit-breaker overhead with 5-run CIs):
run \texttt{python experiments/breaker\_overhead\_multirun.py}.
Output is \texttt{results/overhead\_percentiles/breaker\_multirun.json}.
\item Table~\ref{tab:scalability}: run
\texttt{python experiments/scalability\_stress\_test.py}. Output is
written to \texttt{results/scalability/}.
\item Table~\ref{tab:fdr}, Table~\ref{tab:fdrmatrix}, and
Table~\ref{tab:fdrframework}: run
\texttt{python benchmarks/run\_benchmarks.py}. Output is written to
\texttt{benchmarks/results\_full.tsv}.
\item Multi-agent topology study (\S\ref{sec:realllm}): run
\texttt{python experiments/multi\_agent\_topology\_cli.py}. Output is
written to \texttt{results/multi\_agent\_topology\_cli/}.
\end{itemize}

The Python environment is pinned in \texttt{requirements.lock};
critical pins are \texttt{opentelemetry-api/sdk 1.27.x},
\texttt{Python 3.12}, and the framework versions named in the
adapter source.

\section{Related Work}
\label{sec:related}

\paragraph{Distributed tracing.} Dapper~\cite{dapper} is the canonical
large-scale tracing system and sets the reviewer's mental model.
Canopy~\cite{canopy} introduced schema-based end-to-end performance
tracing at Facebook. Pivot Tracing~\cite{pivottracing} added dynamic
instrumentation and causal aggregation; its happens-before joins are
evaluated post-hoc over collected traces, whereas \sysname's circuit
breaker fires in-process during the same trace as a
\texttt{SpanProcessor} side-effect --- the runtime-control loop is the
mechanical differentiator, not the schema. X-Trace~\cite{xtrace} and
Magpie~\cite{magpie} are the earlier endpoints in the same line.
WAP5~\cite{lighthouse} and \emph{lprof}~\cite{ohmygod_lprof} treat the
inverse problem (reconstruct request flows from logs) and inform our
design's choice to emit explicit spans rather than rely on log
correlation. SEDA~\cite{aegis_seda} and \emph{The Tail at
Scale}~\cite{stout} frame the latency-percentile discipline we apply in
\S\ref{sec:microbench}--\S\ref{sec:scalability}.

\paragraph{OpenTelemetry.} OTel~\cite{otel} is the substrate. The OTel
GenAI semantic conventions~\cite{otel_genai} are the closest prior
art; \sysname{} extends them additively and is the first system, to
our knowledge, to layer a runtime-control SpanProcessor on top.
Jaeger~\cite{jaeger} and Prometheus~\cite{prometheus} are
the downstream observability tooling \sysname{} interoperates with via
OTLP.

\paragraph{Agent observability.} Four contemporaneous systems define
the closest competitive landscape. \textbf{OpenInference}
\cite{openinference} ships an attribute schema (\texttt{openinference.span.kind}
distinguishing \texttt{CHAIN}, \texttt{LLM}, \texttt{RETRIEVER},
\texttt{TOOL}, \texttt{AGENT}, \texttt{RERANKER}) plus auto-instrumentation packages
keyed to a closed set of providers; it does not emit
\texttt{PLANNING}, \texttt{REASONING}, \texttt{DELEGATION},
\texttt{GUARD\_RAIL}, or \texttt{MEMORY}, and ships no runtime-control
mechanism. \textbf{OpenLLMetry}~\cite{openllmetry} wraps a different
schema with similar coverage and similar omissions.
\textbf{LangSmith}~\cite{langsmith} is a managed back-end with a
proprietary vendor-defined trace model; the wire format is closed.
\textbf{Langfuse}~\cite{langfuse} is an open back-end whose trace
model centres on input/output capture and per-call cost; orchestration
phases are not first-class. \textbf{AgentOps}~\cite{agentops} is a
conceptual taxonomy without a runnable implementation we can benchmark
against. None of the five exposes a SpanProcessor that consumes the
live span stream and enforces in-process policies; this is the gap our
\texttt{AgentCircuitBreaker} closes.

\paragraph{Agent failure taxonomies.}
MAST~\cite{mast}, AgentDebug~\cite{agentdebug}, and
Aegis~\cite{aegis} characterize the failure space we instrument
against. Our 14-class fault catalogue is informed by but does not
duplicate any of theirs; we chose classes for which a span-level
signal is constructible, which is a stricter criterion than what those
taxonomies require. None of MAST, AgentDebug, or Aegis ships a
runnable detector against a public schema we could benchmark
head-to-head; our four-baseline comparison (vanilla OTel, OTel GenAI
semconv, OpenInference, OpenLLMetry) is the set of \emph{runnable}
schemas with public attribute conventions sufficient to instantiate the
matcher.

\paragraph{Safety rails.} GuardAgent~\cite{guardagent} and NeMo
Guardrails~\cite{nemo_guardrails} run upstream of the model and
constrain inputs/outputs. Our circuit breaker complements rather than
substitutes for these: GuardAgent enforces semantic constraints on a
single call; the breaker enforces orchestration-level constraints
(cost, loop, cycle) across a trace.

\paragraph{Circuit breaking.}
The circuit-breaker pattern~\cite{circuit_breaker_pattern} is
established in distributed-systems engineering. Our contribution is
not the pattern but its wiring to a tracing pipeline:
\texttt{AgentCircuitBreaker} is a \texttt{SpanProcessor}, which makes
it composable with any OTel-instrumented application at no protocol
cost.

\section{Conclusion}
\label{sec:conclusion}

LLM agent systems fail in ways the existing tracing stack does not
see. Closing the gap requires a span vocabulary that names the
orchestration phases, an adapter layer that absorbs the
heterogeneous-framework binding-strategy space, and a runtime
mechanism that turns the same span stream into enforceable policy.
\sysname{} delivers all three, with measured per-span overhead of
$\sim$11.7\,\textmu{}s p50, a four-policy circuit breaker that costs
+2.16\,\textmu{}s p50 (95\% CI [+2.06, +2.25] across 5 independent
runs), 19{,}071 spans/s under 100-thread concurrency, and an
end-to-end fault-detection rate that ranges from 0.612 across
heterogeneous adapters to 1.000 on a conformance-complete adapter,
against a 0.429 ceiling for the closest production baselines.

\balance

\bibliographystyle{ACM-Reference-Format}
\bibliography{refs}

\end{document}
